\section{Fibonacci syntax and finite Zeckendorf sections}
\label{sec:golden-syntax}

Throughout the paper we use the Fibonacci convention
\[
F_1=F_2=1,\qquad F_{k+1}=F_k+F_{k-1}\quad(k\ge2).
\]
Thus \(1,1,2,3,5,8,\ldots\) is the Fibonacci sequence.

A binary word \(z=z_1\cdots z_m\in\{0,1\}^m\) is called \emph{Zeckendorf-legal} if it contains no adjacent ones:
\[
z_kz_{k+1}\ne 11\qquad(1\le k<m).
\]
The stable type space at resolution \(m\) is
\[
X_m:=\{z\in\{0,1\}^m:z_kz_{k+1}\ne 11\text{ for all }1\le k<m\}.
\]

For \(z\in\{0,1\}^m\), define its shifted Fibonacci value by
\[
V_m(z):=\sum_{k=1}^{m}z_kF_{k+1}.
\]
The shift by one index is chosen so that legal words of length \(m\) represent exactly the interval
\[
\{0,1,\ldots,F_{m+2}-1\}.
\]

\begin{lemma}[Finite Zeckendorf section]
\label{lem:finite-zeckendorf-section}
For every \(m\ge1\), the map
\[
V_m:X_m\to\{0,1,\ldots,F_{m+2}-1\},
\qquad
V_m(z)=\sum_{k=1}^{m}z_kF_{k+1},
\]
is a bijection.  In particular,
\[
|X_m|=F_{m+2}.
\]
\end{lemma}

\begin{proof}
We argue by induction on \(m\).  For \(m=1\), \(X_1=\{0,1\}\), and
\[
V_1(0)=0,\qquad V_1(1)=F_2=1.
\]
Thus \(V_1(X_1)=\{0,1\}=\{0,\ldots,F_3-1\}\).

Assume the claim for all smaller lengths.  Partition \(X_m\) according to the last digit.  If \(z_m=0\), then the prefix \(z_1\cdots z_{m-1}\) is an arbitrary element of \(X_{m-1}\).  By induction, the possible values are
\[
0,1,\ldots,F_{m+1}-1.
\]
If \(z_m=1\), legality forces \(z_{m-1}=0\).  The first \(m-2\) digits form an arbitrary word in \(X_{m-2}\).  Hence the value is
\[
F_{m+1}+r,\qquad 0\le r<F_m.
\]
This case gives precisely
\[
F_{m+1},F_{m+1}+1,\ldots,F_{m+1}+F_m-1.
\]
Since \(F_{m+1}+F_m=F_{m+2}\), the two cases cover \(\{0,1,\ldots,F_{m+2}-1\}\) without overlap.  Therefore \(V_m\) is bijective.
\end{proof}

\begin{definition}[Finite Zeckendorf section]
\label{def:finite-zeckendorf-section}
The inverse of \(V_m\) is denoted
\[
s_m:\{0,1,\ldots,F_{m+2}-1\}\to X_m.
\]
We call \(s_m\) the finite Zeckendorf section.
\end{definition}

\begin{remark}[Role of \(\varphi\)]
The asymptotic growth rate of \(|X_m|\) is \(\varphi=(1+\sqrt5)/2\).  The role of \(\varphi\) in this paper is grammatical: it determines the legal syntax of stable representatives and the normalization section.  No physical interpretation is assumed.
\end{remark}
